\documentclass[10pt,a4paper]{amsart}
\usepackage[margin=19mm]{geometry}
\usepackage{amsmath,amssymb,mathtools}
\usepackage[scr=rsfs,cal=euler]{mathalfa}
\usepackage{xfrac}
\usepackage[unicode,hidelinks,bookmarks=false]{hyperref}
\newcommand{\ie}{i.e.\ }
\newcommand{\cf}{cf.\ }
\newcommand{\resp}{respectively\ }
\newcommand{\Subsection}{\subsection}
\providecommand{\nobreakdash}{\nobreak\hbox{-}}
\setlength{\emergencystretch}{2em}
\multlinegap0pt
\usepackage{mdframed}
\usepackage{mdframed}
\usepackage{mdframed}
\usepackage{mdframed}
\usepackage{amssymb}
\usepackage{mdframed}
\usepackage{bm}
\usepackage{mathtools}
\usepackage{xcolor}
\newtheorem{lemma}{Lemma}
\newcommand{\Brace}[1]{\left\{#1\right\}}
\newcommand{\Brack}[1]{\left[#1\right]}
\newcommand{\Par}[1]{\left(#1\right)}
\newcommand{\R}{\mathbb{R}}
\newcommand{\bbE}{{\mathbb{E}}}
\newcommand{\defeq}{:=}
\newcommand{\eps}{\epsilon}
\newcommand{\sech}{\operatorname{sech}}
\title{High-Magnetization Sampling at Low Temperatures: \\ Ising Models and Bayesian Sparse Linear Regression}
\author{Syamantak Kumar, Purnamrita Sarkar, Kevin Tian, Yusong Zhu}
\date{Source: arXiv:2609.08873v1 (CC BY 4.0)}
\begin{document}
\maketitle

\noindent\emph{Source and licence.} High-Magnetization Sampling at Low Temperatures: \\ Ising Models and Bayesian Sparse Linear Regression. Version arXiv:2609.08873v1 (CC BY 4.0), distributed by arXiv under the Creative Commons Attribution 4.0 International licence, http://creativecommons.org/licenses/by/4.0/, as recorded in the arXiv licence metadata for this exact version and shown on the abstract page https://arxiv.org/abs/2609.08873v1. Full licence notice: \texttt{licenses/arxiv-2609.08873-CC-BY-4.0.txt}.\par\smallskip

\noindent\textit{Editorial scope.} Counted unit: the article's lemma lem:at\_line (source0.tex lines 2609--2623). The article's complete original proof of it is delivered with it (source0.tex lines 2625--2710, one \texttt{proof} environment of 86 lines). No mathematics is written, repaired or re-derived: the statement and the proof are the article's own lines, in the article's own notation, and no retained line is substituted or re-flowed.  Retained uncounted as the source-local context the proof needs: lines 2505--2547; lines 2552--2581.  Nothing was omitted from any retained span, so the package carries no omission record.  No retained line contains a citation, so the wrapper carries no bibliography and no bibliography entry.  No retained line contains a figure environment, an image-inclusion command or drawing code, so no image is delivered and the package declares no asset dependency.  Editorial formatting only, in addition: an AMS \texttt{amsart} preamble that restates the article's own declarations and macros used by the retained lines, namely the environments \texttt{lemma} on separate counters, together with the standard packages those lines need.  The retained environments print in this document as Lemma 1 in order of appearance, whereas the article prints the same items with the numbering of its own full document.  The complete native source of the article as distributed in the arXiv e-print for this exact version is bundled unchanged as \texttt{sources/209758/source0.tex}.
\begin{lemma}
    \label{lem:sech_laplace}
    Let \( \phi(t) \defeq \frac{1}{\sqrt{2\pi}} \exp(-\frac{t^2}{2})\) be the Gaussian density, and let $p\in\Brace{2,4}$. Suppose that there are positive $\{q_\beta, h_\beta\}_{\beta > 0}$ satisfying, as $\beta \to \infty$,
    $q_\beta\to1$, $b_\beta \defeq \frac{h_\beta} \beta \to\infty$, and
    $\frac{b_\beta} \beta\to0$. Then
    \begin{equation}
        \label{eq:sech_laplace}
        \bbE\Brack{
            \operatorname{sech}^p
            \Par{
                \beta\sqrt{q_\beta}Z+\beta b_\beta
            }
        }
        =
        \frac{
            I_p
        }{
            \beta\sqrt{q_\beta}
        }
        \phi\Par{
            \frac{
                b_\beta
            }{
                \sqrt{q_\beta}
            }
        }
        \Par{1+o(1)},
    \end{equation}
    where
    \[
        I_2
        =
        \int_{\R}\operatorname{sech}^2(u)\,\mathrm du
        =
        2,
        \qquad
        I_4
        =
        \int_{\R}\operatorname{sech}^4(u)\,\mathrm du
        =
        \frac{4}{3}.
    \]
\end{lemma}

\begin{equation}\label{eq:cov_integral}
    \bbE\Brack{
        \operatorname{sech}^p
        \Par{
            \beta\sqrt{q_\beta}Z+\beta b_\beta
        }
    }
    =
    \frac{
        \phi\Par{b_\beta/\sqrt{q_\beta}}
    }{
        \beta\sqrt{q_\beta}
    }
    \int_{\R}
    \operatorname{sech}^p(u)
    \exp\Par{
        \frac{
            b_\beta u
        }{
            \beta q_\beta
        }
        -
        \frac{
            u^2
        }{
            2\beta^2q_\beta
        }
    }
    \mathrm du.
\end{equation}

\begin{lemma}
    \label{lem:at_line}
    As $\beta\to\infty$,
    \[
        h_{\rm AT}(\beta)
        =
        \sqrt{2}\,\beta\sqrt{\log\beta}
        \Par{1+O\Par{\frac{1}{\log\beta}}},
        \quad
        h_{\rm wAT}(\beta)
        =
        \sqrt{2}\,\beta\sqrt{\log\beta}
        \Par{1+O\Par{\frac{1}{\log\beta}}}.
    \]
\end{lemma}

\begin{proof}
We begin by verifying the conditions of Lemma~\ref{lem:sech_laplace}. For $\star\in\Brace{\mathrm{AT},\mathrm{wAT}}$, write
\[
    q_\star\defeq q_\star(\beta),
    \quad
    b_\star\defeq \frac{h_\star(\beta)}{\beta}.
\]
We also drop the index $\beta$ from these two sequences for simplicity.
First, along the weak AT boundary,
$1-q_{\rm wAT}=\beta^{-2}$, so $q_{\rm wAT}\to1$. Similarly, along the AT boundary,
\[
    \bbE\Brack{
        \operatorname{sech}^4
        \Par{
            \beta\sqrt {q_{\rm AT}} Z+h_{\rm AT}(\beta)
        }
    }
    =
    \frac{1}{\beta^2}.
\]
Since
$1-q_{\rm AT}=\bbE[\operatorname{sech}^2(\beta\sqrt {q_{\rm AT}} Z+h_{\rm AT}(\beta))]$,
Cauchy-Schwarz gives $1-q_{\rm AT}\le\beta^{-1}$, so
$q_{\rm AT}\to1$.

Next, for either value of $\star$, the boundary equation also implies
$b_\star \to\infty$. Indeed, if $b_\star$
were bounded along a subsequence, then restricting the integral \eqref{eq:cov_integral} to
$u\in[-1,1]$ would already give a lower bound of order $\beta^{-1}$, because $q_\star \to 1$, $b_\star$ is bounded, and $\sech^p(u) = \Omega(1)$ for $u \in [-1, 1]$. This would
contradict the boundary equations, which say that \eqref{eq:cov_integral} evaluates to $\beta^{-2}$.


Finally, we claim that $\frac{b_\star}{\beta} \to0$.  Otherwise, $b_\star\ge\eps\beta$ along a subsequence
for some $\eps>0$. Now, split the integral \eqref{eq:cov_integral} along the events $Z \ge -b_\star / 2\sqrt{q_\star}$ or $Z \le -b_\star / 2\sqrt{q_\star}$. In the former region, the change of variables gives $u \ge \frac{\eps\beta^2}{2}$, so the corresponding integral is $\exp(-\Omega(\beta^2))$. Similarly, the latter region has a probability bounded by $\exp(-\Omega(\beta^2))$, and $\sech$ is pointwise bounded. Thus, the entire integral is $\exp(-\Omega(\beta^2))$, again contradicting that it equals $\beta^{-2}$ by definition.

Lemma~\ref{lem:sech_laplace} therefore applies and we obtain that 
\[
    \phi\Par{
        \frac{
            b_{\rm AT}
        }{
            \sqrt {q_{\rm AT}}
        }
    }
    =
    \frac{
        3\sqrt {q_{\rm AT}}
    }{
        4\beta
    }
    \Par{1+o(1)}, 
    \quad 
    \phi\Par{
        \frac{
            b_{\rm wAT}
        }{
            \sqrt {q_{\rm wAT}}
        }
    }
    =
    \frac{
        \sqrt {q_{\rm wAT}}
    }{
        2\beta
    }
    \Par{1+o(1)}.
\]
Expanding with the definition of $\phi$, and then taking logarithms, then gives for
$\star\in\Brace{\mathrm{AT},\mathrm{wAT}}$,
\[
    \frac{b_\star^2}{2q_\star}
    =
    \log\beta+O(1) \implies b_\star^2 = 2q_\star \log\beta + O(1) = 2\log\beta - 2(1-q_\star)\log\beta + O(1).
\]
Since $1-q_{\rm AT}\le\beta^{-1}$ and
$1-q_{\rm wAT}=\beta^{-2}$, the desired claims follow:
\[
    b_\star^2
    =
    2\log\beta+O(1) \implies
    b_\star
    =
    \sqrt{2\log\beta}
    \Par{1+O\Par{\frac{1}{\log\beta}}}.
\]
\end{proof}

\end{document}
